\begin{theorem}[全息 Kraft--Gödel 侧信息长度下界]\label{thm:hypercube-kraft-godel-sideinfo-lower-bound}
\leanverified{paper\_hypercube\_kraft\_godel\_sideinfo\_lower\_bound}
令 \(f:\Omega_n=\{0,1\}^n\twoheadrightarrow X\) 为满射粗粒化，并沿用界面面积 \(A(f)\) 与规范体积密度 \(g(f)\) 的记号（定义 \ref{def:hypercube-coarsegraining-area-gauge-volume}）。
固定字母表 \(\Sigma\)，\(\card{\Sigma}=M\ge 2\)。
设存在映射
\[
\Phi:\Omega_n\longrightarrow X\times \Sigma^\ast
\]
满足 \(\Phi(\omega)=\bigl(f(\omega),\Phi_2(\omega)\bigr)\) 且 \(\Phi\) 在 \(\Omega_n\) 上单射，并且对每个 \(x\in X\)，集合 \(\{\Phi_2(\omega):\omega\in f^{-1}(x)\}\subset \Sigma^\ast\) 为前缀码。
定义平均侧信息长度
\[
L(\Phi):=2^{-n}\sum_{\omega\in\Omega_n}\bigl|\Phi_2(\omega)\bigr|.
\]
则有下界
\[
\boxed{\;
L(\Phi)\ \ge\ \frac{g(f)+1-\card{X}/2^n}{\log M}
\ \ge\ \frac{n\log 2-\frac{\log 2}{2^{n-1}}A(f)}{\log M}.\;}
\]
特别地，当 \(M=2\) 时，
\[
L(\Phi)\ \ge\ n-\frac{A(f)}{2^{n-1}}.
\]
\end{theorem}
\begin{proof}
对固定 \(x\in X\) 记纤维大小 \(d_x:=\card{f^{-1}(x)}\)，并将纤维上的码字长度记为 \(\ell_{x,1},\dots,\ell_{x,d_x}\)。
由于 \(\{\Phi_2(\omega):\omega\in f^{-1}(x)\}\) 为 \(M\)-元前缀码，Kraft 不等式给出
\[
\sum_{j=1}^{d_x}M^{-\ell_{x,j}}\le 1.
\]
注意 \(t\mapsto M^{-t}=\exp(-t\log M)\) 为凸函数，Jensen 不等式得到
\[
\frac1{d_x}\sum_{j=1}^{d_x}M^{-\ell_{x,j}}
\ \ge\
M^{-\frac1{d_x}\sum_{j=1}^{d_x}\ell_{x,j}}
=M^{-L_x},
\qquad
L_x:=\frac1{d_x}\sum_{j=1}^{d_x}\ell_{x,j}.
\]
从而
\[
1\ \ge\ \sum_{j=1}^{d_x}M^{-\ell_{x,j}}\ \ge\ d_x\,M^{-L_x},
\qquad\text{即}\qquad
L_x\ge \log_M d_x=\frac{\log d_x}{\log M}.
\]
令 \(p_x:=d_x/2^n\)，并对 \(x\) 加权求和得
\[
L(\Phi)=\sum_{x\in X}p_x L_x\ \ge\ \frac{1}{\log M}\sum_{x\in X}p_x\log d_x.
\]
另一方面，对整数 \(m\ge 1\) 有基本上界 \(\log(m!)\le m\log m-m+1\)，故
\[
g(f)=2^{-n}\sum_{x\in X}\log(d_x!)
\le
2^{-n}\sum_{x\in X}\bigl(d_x\log d_x-d_x+1\bigr)
=\sum_{x\in X}p_x\log d_x-1+\frac{\card{X}}{2^n},
\]
即 \(\sum_x p_x\log d_x\ge g(f)+1-\card{X}/2^n\)。
合并得到第一条不等式。
最后由命题 \ref{prop:hypercube-discrete-holographic-stokes-gauge-area}，
\[
g(f)\ \ge\ (n\log 2-1)-2\log 2\cdot\frac{A(f)}{2^n}+\frac{\card{X}}{2^n},
\]
从而
\[
g(f)+1-\frac{\card{X}}{2^n}\ \ge\ n\log 2-\frac{\log 2}{2^{n-1}}A(f),
\]
代回得到第二条不等式。证毕。
\end{proof}

\endinput
